% !TeX program = xelatex
% ============================================================
%  3. МАТЕРИАЛЫ (источники)
%
%  Статьи, книги, ссылки. У каждой записи четыре блока:
%    • библиографическое описание и ссылка/DOI;
%    • краткий нейросетевой обзор;
%    • зачем это нужно в работе.
%
% ============================================================

\section{Материалы}
\label{sec:materials}

\subsection{Бахадуровская эффективность}

\subsubsection*{1. Bahadur R.~R. A note on the asymptotic efficiency of statistical tests}
\textit{The Annals of Mathematical Statistics}, 1960, vol.~31, no.~2, pp.~499--503.
\textbf{Обзор.} Вводит сравнение тестов по скорости стремления достигнутого уровня
значимости к нулю; предел $-\frac{2}{n}\log L_n$ называется точным наклоном, отношение
наклонов~--- бахадуровской эффективностью.\\
\textbf{Зачем в работе.} Отправная точка подзадачи 2.1: нужно показать, что для статистик с
невырожденным предельным распределением эта конструкция вырождается, и обосновать переход к
приближённому наклону.

\subsubsection*{2. Bahadur R.~R. Stochastic comparison of tests}
\textit{The Annals of Mathematical Statistics}, 1960, vol.~31, no.~2, pp.~276--295.
DOI: 10.1214/aoms/1177705894.\\
\textbf{Обзор.} Общая схема стохастического сравнения тестов: тест эффективнее, если его
статистика быстрее уходит от критической области.\\
\textbf{Зачем в работе.} Терминологическая рамка раздела «Численный метод». Библиографические
данные подтверждены через Crossref.

\subsubsection*{3. Bahadur R.~R. Rates of convergence of estimates and test statistics}
\textit{The Annals of Mathematical Statistics}, 1967, vol.~38, no.~2, pp.~303--324.\\
\textbf{Обзор.} Связывает скорость сходимости оценок и статистик с асимптотической
эффективностью; скорость убывания хвостовых вероятностей~--- более тонкая характеристика,
чем асимптотическая дисперсия.\\
\textbf{Зачем в работе.} Аргумент во введении: почему сравнения по мощности на конечной
выборке недостаточно.

\subsubsection*{4. Nikitin Ya.~Yu. Asymptotic Efficiency of Nonparametric Tests}
Cambridge University Press, 1995. ISBN 0-521-47029-3.\\
\textbf{Обзор.} Монография: наклон Бахадура, функции больших уклонений, локальная
эффективность, эффективность непараметрических критериев согласия и независимости.\\
\textbf{Зачем в работе.} Основная теоретическая база задачи~2; источник формулировок условий
регулярности, нарушаемых при составной нулевой гипотезе.

\subsubsection*{5. Milošević B., Nikitin Ya.~Yu., Obradović M. Bahadur efficiency of EDF
based normality tests when parameters are estimated}
arXiv:2106.07437, 2021. \texttt{https://arxiv.org/abs/2106.07437}\\
\textbf{Обзор.} Для EDF-критериев нормальности при неизвестных параметрах, где
функции больших уклонений неизвестны, применяется приближённый наклон: если
$\log(1-F(t)) = -\frac{a_T t^2}{2}(1+o(1))$ и $T_n/\sqrt{n} \to b_T(\theta) > 0$, то
$c_T^{*}(\theta) = a_T b_T^2(\theta)$; эталон~--- тест отношения правдоподобий с локальным
точным наклоном $2K(\theta)$.\\
\textbf{Зачем в работе.} \textbf{Главный методологический источник.} Даёт ровно ту схему,
которую нужно реализовать численно: $a_T$ из хвоста предельного распределения (в PySATL оно
уже моделируется и сохраняется в БД), $b_T(\theta)$ из выборок по альтернативе (их уже
генерирует шаг генерации). Определяет содержание подзадач 2.1--2.3 и формат матрицы
эффективностей в подзадаче 5.3.

\subsubsection*{6. Obradović M. On asymptotic efficiency of goodness of fit tests for Pareto
distribution based on characterizations}
\textit{Filomat}, 2015, vol.~29, no.~10, pp.~2311--2324.\\
\textbf{Обзор.} Рассчитаны бахадуровские эффективности критериев согласия распределению
Парето против различных альтернатив; для каждого критерия найдены локально оптимальные
альтернативы.\\
\textbf{Зачем в работе.} Образец результата подзадачи 5.3: таблица «критерий $\times$
альтернатива $\to$ эффективность» плюс локально оптимальные альтернативы.

\subsubsection*{7. Durio A., Nikitin Ya.~Yu. Local Bahadur efficiency of some goodness-of-fit
tests under skew alternatives}
\textit{Journal of Statistical Planning and Inference}, 2003, vol.~115, no.~2,
pp.~171--179. DOI: 10.1016/S0378-3758(02)00155-6.\\
\textbf{Обзор.} Локальная бахадуровская эффективность критериев согласия при
асимметричных альтернативах.\\
\textbf{Зачем в работе.} Источник семейства асимметричных альтернатив для матрицы; важен для
подзадачи 5.2, где локальный режим наиболее неустойчив численно.

\subsubsection*{8. Akritas M.~G., Kourouklis S. Local Bahadur efficiency of score tests}
\textit{Journal of Statistical Planning and Inference}, 1988, vol.~19, no.~1,
pp.~187--199. DOI: 10.1016/0378-3758(88)90072-9.\\
\textbf{Обзор.} Локальный наклон скор-тестов; соотношение локальной и нелокальной
эффективности.\\
\textbf{Зачем в работе.} Обоснование того, что локальная и точная эффективности могут
совпадать,~--- важно для интерпретации расхождений с аналитическими эталонами
(подзадача 5.1).

\subsubsection*{9. He X., Shao Q.-M. Bahadur efficiency and robustness of studentized score
tests}
\textit{Annals of the Institute of Statistical Mathematics}, 1996, vol.~48, no.~2,
pp.~295--314. DOI: 10.1007/BF00054792.\\
\textbf{Обзор.} Выведены точные наклоны студентизированных скор-тестов; отмечена прямая
связь бахадуровской эффективности с достигнутым уровнем значимости.\\
\textbf{Зачем в работе.} Пример полного вывода наклона, на фоне которого видно, какие шаги
заменяются численной оценкой; подтверждает корректность оценки через достигнутые уровни.

\subsubsection*{10. Nikitin Ya.~Yu. Bahadur asymptotic efficiency of integral tests for
symmetry}
\textit{Journal of Soviet Mathematics}, 1984, vol.~27. DOI: 10.1007/BF01843558.\\
\textbf{Обзор.} Наклон для интегральных критериев симметрии выражается через ядро
статистики.\\
\textbf{Зачем в работе.} Часть критериев в \texttt{pysatl-criterion} относится к тому же
классу; полезно при выборе семейства для целевого прогона (задача 5.3).


